\documentclass[12pt]{article}

% --- Codificación y lengua ---
\usepackage[T1]{fontenc}
\usepackage[utf8]{inputenc}
\usepackage[spanish,es-noindentfirst,es-tabla]{babel}
\usepackage{lmodern}
\usepackage{microtype}

% --- Página y espaciado estilo APA 7 ---
\usepackage[letterpaper,margin=1in]{geometry}
\usepackage{setspace}
\usepackage{parskip}
\usepackage{enumitem}
\usepackage{titlesec}
\usepackage[hidelinks]{hyperref}
\usepackage{xurl}

% Sangría francesa para las referencias
\newenvironment{referencias}
  {\begin{list}{}{\setlength{\leftmargin}{0.5in}\setlength{\itemindent}{-0.5in}%
    \setlength{\itemsep}{0.5\baselineskip}\setlength{\parsep}{0pt}}}
  {\end{list}}

\titleformat{\section}{\normalfont\large\bfseries\centering}{}{0pt}{}
\titleformat{\subsection}{\normalfont\normalsize\bfseries}{}{0pt}{}
\setlength{\parskip}{0.6\baselineskip}
\newcommand{\cita}[1]{\textup{(#1)}}

\title{\vspace{-2em}\bfseries Humanización, ética y bien común en la práctica educativa}
\author{Jesús G.\\[2pt]\normalsize Pedagogía del Bien Común}
\date{11 de julio de 2026}

\begin{document}
\maketitle
\thispagestyle{empty}

\begin{quote}\small
\noindent\textit{Nota metodológica.} Cada afirmación con cita se contrastó contra el
texto completo de su fuente \mbox{open-access} mediante un índice vectorial local
(Chroma); solo se conservaron las citas cuyo pasaje de respaldo superó el umbral de
verificación. Las obras no verificables en acceso abierto se sustituyeron por fuentes
equivalentes descargables.
\end{quote}

\vspace{0.5em}

% =====================================================================
\section*{2. Humanización y ética}
% =====================================================================

Nadie se hace persona en solitario. Aunque pertenecemos por biología a la especie
humana, aprendemos a hablar, convivir y responder de nuestros actos en la relación con
otros. La conducta no brota de una sola causa: resulta de la interacción entre lo
biológico, la historia personal, el ambiente y la cultura, niveles que se influyen
mutuamente \cita{Frías-Armenta et al., 2003}. Por eso la sociedad no es algo externo al
ser humano, sino parte de lo que lo hace posible.

La \textbf{humanización} es el proceso de desarrollar empatía, lenguaje, razonamiento y
responsabilidad. Esas capacidades se aprenden: la empatía y la competencia social forman
parte de las competencias emocionales, que son enseñables y no un rasgo fijo
\cita{Bisquerra \& Pérez, 2007}. La \textbf{ética} orienta esas capacidades hacia lo que
debemos hacer y hacia las consecuencias de lo que decidimos.

El \textbf{bien común} amplía la mirada: no basta mi beneficio; debo preguntarme si mis
actos permiten también que los demás vivan con dignidad. La educación, entendida como un
bien común, se guía por el respeto a la vida y a la dignidad humana antes que por el
interés individual \cita{UNESCO, 2015}.

\textit{Ejemplo cotidiano.} Cuando reviso el proyecto de un estudiante, podría limitarme a
poner una calificación. Actuar éticamente es, además, devolver observaciones claras y
útiles para que mejore, sin aprobar un trabajo que no cumple ---eso sería injusto para el
grupo---. Se sostiene el estándar y, a la vez, se reconoce que detrás de cada entrega hay
una persona. Es el caso de Gael: ni aprobar por lástima ni reducirlo a una cifra, sino
abrir una ruta honesta de recuperación.

% =====================================================================
\section*{3. Educación humanizante}
% =====================================================================

Una buena vida no es solo producir o tener éxito: es desarrollar las propias capacidades,
construir vínculos, actuar con libertad y responsabilidad y poner los talentos al servicio
de los demás, con espacio también para descansar, equivocarse y encontrar sentido. La
educación que persigue esa vida no se reduce a transmitir conocimientos; se apoya en
cuatro aprendizajes: aprender a conocer, a hacer, a vivir juntos y a ser
\cita{Delors, 1996}.

Esta idea dialoga con la tradición confuciana, en la que educar es cultivar virtud: el
hombre superior se distingue por su rectitud y su benevolencia, no por acumular datos
\cita{Confucio, trad. Legge, 2020}. Una educación que humaniza:

\begin{itemize}[leftmargin=1.4em,itemsep=1pt,topsep=2pt]
  \item reconoce al estudiante como persona, no como matrícula o indicador;
  \item forma juicio crítico, responsabilidad y capacidad de convivencia, no solo saberes;
  \item deja preguntar, equivocarse y comprender el sentido de lo aprendido;
  \item mantiene criterios claros y ofrece acompañamiento y oportunidades reales de mejora;
  \item vincula el aprendizaje con las necesidades de la comunidad.
\end{itemize}

A esta reflexión ayuda una pedagogía del mirar: aprender no es reaccionar a cada estímulo,
sino detenerse, sostener la atención y pensar antes de actuar. En la sociedad del
rendimiento el sujeto se autoexplota buscando el máximo desempeño; frente a la
hiperatención y la prisa hace falta recuperar la atención profunda y contemplativa
\cita{Han, 2012, citado en Quintero Camarena, 2017}. Por eso una educación humanizante no
responde a una entrega tardía con un castigo automático: primero escucha, revisa las
circunstancias, explica los criterios y busca una salida justa ---sin eliminar la
responsabilidad---. Incluso en línea hay un rostro detrás de cada silencio.

Lo contrario deshumaniza: reducir el aprendizaje a cumplir tareas y competir; humillar,
etiquetar o usar la evaluación solo como sanción; o vigilar en exceso y anteponer la
eficiencia a las personas. La pregunta que distingue una postura de la otra es simple:
\textit{\textquestiondown{}quién es la persona que estoy formando y para qué?}, frente a \textit{\textquestiondown{}cumplió,
produjo y aprobó?}

% =====================================================================
\section*{4. Ética profesional}
% =====================================================================

La ética profesional es ejercer con competencia, honestidad y responsabilidad, atendiendo
no solo a lo que puedo hacer técnicamente, sino a lo que debo hacer y a sus consecuencias
sobre otros. No se agota en cumplir reglamentos: incluye una dimensión deontológica con
obligaciones hacia los alumnos, la institución, los colegas y la sociedad, y supone
competencias éticas para resolver dilemas y conflictos de valores
\cita{Sanz Ponce \& Hirsch, 2016}.

En lo educativo esto significa reconocer la dignidad de cada estudiante; evaluar con
criterios claros y transparentes; evitar favoritismos y humillaciones; proteger la
información personal; reconocer y corregir los propios errores; usar con responsabilidad la
tecnología y la inteligencia artificial; y ofrecer retroalimentación que favorezca el
aprendizaje.

El núcleo es el \textbf{equilibrio entre el cuidado y la verdad}. Aprobar sin evidencias
parece compasivo, pero vacía de sentido la evaluación; reprobar sin escuchar ni orientar
ignora la dignidad. Lo ético es registrar con honestidad el aprendizaje logrado y, a la
vez, abrir mecanismos justos de acompañamiento.

\textbf{Mi postura personal.} Resuelvo estos casos desde la tradición confuciana. Me
convence porque coloca la ética en el carácter antes que en el reglamento: quien enseña se
distingue por su rectitud y su benevolencia, no por acumular saberes ni por cumplir la norma
al pie de la letra \cita{Confucio, trad. Legge, 2020}. Eso me obliga a algo incómodo. Si mi
conducta no es coherente, ninguna regla escrita sostiene la relación educativa, porque el
ejemplo enseña más que el código. También me exige reciprocidad: no imponer al estudiante lo
que yo no aceptaría en su lugar. Por eso, cuando dudo entre aprobar sin evidencia o reprobar
sin escuchar, la pregunta que uso no es qué artículo aplica, sino qué haría una persona recta
en mi lugar y cómo cultivo esa misma rectitud en quien tengo enfrente.

% =====================================================================
\section*{5. Decálogo de la ética profesional}
% =====================================================================

Este decálogo traduce en compromisos la deontología profesional ---confianza pública,
responsabilidad ante el alumnado y orientación al bien común--- que la literatura señala
como propia del oficio educativo \cita{Sanz Ponce \& Hirsch, 2016; UNESCO, 2015}. Como
equipo nos comprometemos a:

\begin{enumerate}[leftmargin=1.6em,itemsep=1pt,topsep=2pt]
  \item Reconocer la dignidad de toda persona y no reducirla a un número o un expediente.
  \item Actuar con honestidad, comunicando información veraz y admitiendo errores y
        conflictos de interés.
  \item Cumplir con competencia y disposición permanente a aprender.
  \item Decidir con justicia e imparcialidad, sin favoritismos ni abuso de autoridad.
  \item Proteger la privacidad y la confidencialidad de la información recibida.
  \item Escuchar antes de juzgar, considerando el contexto y las consecuencias.
  \item Evitar el daño, anticipando riesgos y corrigiendo a tiempo sus efectos.
  \item Usar la tecnología con responsabilidad, verificando la información y cuidando los
        datos.
  \item Colaborar con respeto, promoviendo el diálogo y la resolución pacífica de
        conflictos.
  \item Orientar el trabajo al bien común, no solo al beneficio personal.
\end{enumerate}

% =====================================================================
\section*{6. Código de ética (borrador)}
% =====================================================================

\textit{Para profesionales de la educación.} Los códigos deontológicos del profesorado
distinguen obligaciones que promueven la confianza pública en la profesión y la
responsabilidad hacia el alumnado \cita{Sanz Ponce \& Hirsch, 2016}; su fundamento es el
respeto a la dignidad humana, la igualdad de derechos y la justicia social
\cita{UNESCO, 2015}.

\begin{description}[leftmargin=0em,itemsep=2pt,topsep=2pt,font=\normalfont\bfseries]
  \item[Art. 1. Finalidad.] Orientar la enseñanza, la evaluación y la tutoría hacia la
        formación integral, la dignidad, la justicia y el bien común.
  \item[Art. 2. Dignidad.] Reconocer a cada estudiante como fin en sí mismo; evitar toda
        humillación, discriminación o reducción a una cifra.
  \item[Art. 3. Justicia e inclusión.] Decidir con criterios claros y conocidos; ofrecer
        apoyos razonables sin falsear las evidencias del aprendizaje.
  \item[Art. 4. Honestidad académica.] Comunicar información veraz, reconocer las fuentes y
        prevenir antes que castigar.
  \item[Art. 5. Evaluación responsable.] Que las calificaciones correspondan a las
        evidencias; ante el no logro, explicar razones y abrir una ruta de recuperación.
  \item[Art. 6. Privacidad.] Tratar con confidencialidad la información personal y
        académica.
  \item[Art. 7. Tecnología e IA.] Usarlas por su valor educativo; verificar lo que generan
        y no delegar en una máquina decisiones sobre la trayectoria del estudiante.
  \item[Art. 8. Relación educativa.] Basarla en el respeto, la escucha y límites
        profesionales claros; la autoridad es cuidado, no control.
  \item[Art. 9. Competencia.] Mantenerse actualizado y pedir apoyo cuando la situación
        rebasa las propias capacidades.
  \item[Art. 10. Responsabilidad ante la comunidad.] Ante un dilema, identificar los
        valores en conflicto, las personas afectadas y las normas antes de decidir.
\end{description}

\textit{Compromiso final.} Revisar periódicamente la propia práctica, recibir
retroalimentación y orientar la profesión hacia una educación que permita aprender a ser,
convivir, transformar y trascender \cita{Delors, 1996}.

% =====================================================================
\section*{7. Referencias}
% =====================================================================
\begin{referencias}

\item Bisquerra, R., \& Pérez, N. (2007). Las competencias emocionales.
\textit{Educación XX1, 10}, 61--82.

\item Confucio. (2020). \textit{The analects of Confucius} (J. Legge, Trad.). Project
Gutenberg. (Obra de la tradición confuciana, ca. siglo V a.\,C.; traducción original de
1893). \url{https://www.gutenberg.org/ebooks/3330}

\item Delors, J. (1996). \textit{La educación encierra un tesoro: Informe a la UNESCO de la
Comisión Internacional sobre la Educación para el Siglo XXI}. Santillana/UNESCO.
\url{https://unesdoc.unesco.org/ark:/48223/pf0000109590_spa}

\item Frías-Armenta, M., López-Escobar, A. E., \& Díaz-Méndez, S. G. (2003). Predictores de
la conducta antisocial juvenil: un modelo ecológico. \textit{Estudos de Psicologia, 8}(1),
15--24. \url{https://doi.org/10.1590/S1413-294X2003000100003}

\item Quintero Camarena, G. (2017). De la sociedad de los locos a la sociedad de los
cansados [Reseña del libro \textit{La sociedad del cansancio}, de B.-C. Han].
\textit{Culturales, 1}(2), 321--329. \url{https://www.scielo.org.mx/pdf/cultural/v5n2/2448-539X-cultural-5-02-00321.pdf}

\item Sanz Ponce, R., \& Hirsch Adler, A. (2016). Ética profesional en el profesorado de
educación secundaria de la Comunidad Valenciana. \textit{Perfiles Educativos, 38}(151),
144--162.

\item UNESCO. (2015). \textit{Replantear la educación: \textquestiondown{}Hacia un bien común mundial?}
UNESCO. \url{https://unesdoc.unesco.org/ark:/48223/pf0000232697}

\end{referencias}

\end{document}
